\documentclass{../templateWriting/cdwriting}

\newcommand{\authorinfo}{THE FORUM CENTER - NGUYỄN HOÀNG HUY}
\newcommand{\documenttitle}{TÀI LIỆU CHUYÊN ĐỀ WRITING}
\newcommand{\collectiontitle}{Level 3 — W6 (Feb 2-8)}

\begin{document}

\thispagestyle{firstpage}
\makeworkshopheader{\authorinfo}{\documenttitle}{\collectiontitle}
\pagestyle{otherpages}

\workshopprompt{Writing, reading and maths are the three recognized traditional subjects. Computer skill should also be the fourth largest branch. Do you agree or disagree?}

\cdstep{1. Yêu cầu cốt lõi}

Đề bài yêu cầu bạn định vị vai trò của "Kỹ năng máy tính" (Computer skills) ngang hàng với 3 môn học cốt lõi truyền thống (Đọc, Viết, Toán).
\begin{itemize}
\item \cdkv{Dạng câu hỏi}{"Do you agree or disagree?"}
\item \cdkv{Lưu ý quan trọng}{Với dạng câu hỏi trực tiếp này, chiến lược an toàn và dễ đạt điểm cao nhất về tiêu chí Task Response là chọn Strong Opinion (Hoàn toàn đồng ý hoặc Hoàn toàn không đồng ý).}
\begin{itemize}
\item Bạn không nên chọn "Partly agree" (nửa vời) nếu không đủ khả năng ngôn ngữ để xử lý sự phức tạp, vì dễ dẫn đến mâu thuẫn trong lập luận.
\item \cdkv{Chiến lược bài này}{Chọn HOÀN TOÀN ĐỒNG Ý (Totally Agree). Chúng ta sẽ chứng minh rằng trong kỷ nguyên số, mù tin học cũng tệ hại như mù chữ.}
\end{itemize}
\end{itemize}

\cdstep{2. Lỗi thường gặp ở Band 5-6}
\begin{itemize}
\item \cdkv{Lạc đề}{Chỉ liệt kê lợi ích của máy tính (chơi game, giải trí...) mà không chứng minh được tại sao nó nên là "môn học cốt lõi thứ 4" (bắt buộc, nền tảng).}
\item \cdkv{Lập trường không nhất quán}{Mở bài nói đồng ý, nhưng thân bài lại dành quá nhiều đất diễn để nói về tác hại của game online \$\textbackslash\{\}rightarrow\$ Mất điểm Coherence.}
\end{itemize}

\cdsection{PHẦN I: PHÂN TÍCH ĐỀ VÀ TÌM Ý \textit{(BRAINSTORMING)}}

Chúng ta cần tìm các luận điểm sắc bén để thuyết phục giám khảo rằng: Không có tin học = Thất bại trong xã hội hiện đại.

\cdblue{\textbf{Luận điểm chính (Main Ideas) cho phe ĐỒNG Ý:}}

\cdstep{1. Tại sao nó quan trọng đối với sự nghiệp \textit{(Career \& Economy)}?}
\begin{itemize}
\item \cdkv{Analysis}{Hầu hết mọi ngành nghề hiện nay (từ bác sĩ, kỹ sư đến nhân viên văn phòng) đều dùng phần mềm.}
\item \cdkv{Key concept}{"Digital literacy" (Sự hiểu biết về kỹ thuật số) hiện nay là một điều kiện tiên quyết (prerequisite) để có việc làm, giống như việc biết đọc biết viết ngày xưa.}
\item \cdkv{Example}{Kỹ năng xử lý dữ liệu, gửi email, dùng phần mềm quản lý.}
\end{itemize}

\cdstep{2. Tại sao nó quan trọng đối với giáo dục và đời sống \textit{(Education \& Daily Life)}?}
\begin{itemize}
\item \cdkv{Analysis}{Máy tính là công cụ để học các môn khác (tra cứu, làm bài tập). Ngoài ra, các dịch vụ công, ngân hàng, y tế đều đã số hóa.}
\item \cdkv{Key concept}{"The fourth pillar" (Cột trụ thứ 4). Nếu không biết dùng máy tính, con người sẽ bị cô lập khỏi xã hội (socially excluded).}
\item \cdkv{Comparison}{So sánh trực tiếp với Đọc/Viết. Biết đọc để tiếp thu tri thức sách vở \$\textbackslash\{\}rightarrow\$ Biết dùng máy tính để tiếp thu tri thức nhân loại trên Internet.}
\end{itemize}

\cdsection{PHẦN II: CẤU TRÚC BÀI VIẾT \textit{(OUTLINE)}}

\cdstep{1. Mở bài \textit{(Introduction)}:}
\begin{itemize}
\item \cdkv{Hook}{Giới thiệu bối cảnh giáo dục truyền thống (3 môn chính).}
\item \cdkv{Thesis Statement}{Khẳng định dứt khoát quan điểm: Tôi hoàn toàn đồng ý rằng kỹ năng máy tính cần được đưa vào chương trình giảng dạy cốt lõi ngang hàng với Đọc, Viết và Toán.}
\end{itemize}

\cdstep{2. Thân bài 1 \textit{(Professional Necessity)}: Tầm quan trọng trong thị trường lao động.}
\begin{itemize}
\item Máy tính thống trị nơi làm việc (workplace dominance).
\item Thiếu kỹ năng này dẫn đến thất nghiệp hoặc hạn chế cơ hội thăng tiến.
\end{itemize}

\cdstep{3. Thân bài 2 \textit{(Educational \& Social Necessity)}: Tầm quan trọng trong việc tiếp cận tri thức và đời sống.}
\begin{itemize}
\item Máy tính là cổng kết nối tri thức (gateway to information).
\item Hỗ trợ việc học các môn truyền thống hiệu quả hơn.
\item Cần thiết cho các hoạt động đời sống hàng ngày (ngân hàng, giao tiếp).
\end{itemize}

\cdstep{4. Kết bài \textit{(Conclusion)}:}
\begin{itemize}
\item \cdkv{Restate}{Khẳng định lại tầm quan trọng không thể thay thế của kỹ năng máy tính.}
\item \cdkv{Final thought}{Nó không còn là môn phụ (optional subject) mà là kỹ năng sinh tồn (survival skill).}
\end{itemize}

\cdsection{PHẦN III: BÀI MẪU CHI TIẾT VÀ PHÂN TÍCH}

\cdstep{1. Introduction}

"Traditionally, literacy and numeracy—comprising reading, writing, and mathematics—have been regarded as the cornerstones of formal education. However, in the digital era, I firmly believe that computer proficiency has become equally indispensable. Therefore, I agree that it should be officially recognized as the fourth fundamental pillar of the school curriculum."
\begin{itemize}
\item \cdblue{\textbf{Phân tích:}}
\begin{itemize}
\item \cdkv{Paraphrase}{"Writing, reading and maths" -> "Literacy and numeracy".}
\item \cdkv{Vocab}{Cornerstones (nền tảng/đá tảng), Indispensable (không thể thiếu), Proficiency (sự thành thạo).}
\end{itemize}
\end{itemize}

\cdstep{2. Body Paragraph 1: Yếu tố nghề nghiệp \textit{(Employability)}}

"The primary rationale for elevating computer skills to a core subject status is their ubiquity in the modern workforce. In today's knowledge-based economy, digital literacy is no longer an optional asset but a prerequisite for employability across virtually all sectors. From basic administrative tasks to complex data analysis, the ability to operate technological devices is essential. Without this competence, job seekers would face significant disadvantages, similar to being illiterate in the past. Consequently, integrating computer science into the core curriculum ensures that future generations are adequately prepared for the demands of the professional world."
\begin{itemize}
\item \cdblue{\textbf{Phân tích:}}
\begin{itemize}
\item \cdkv{Topic Sentence}{Rõ ràng, tập trung vào "ubiquity in the modern workforce" (sự phổ biến trong lực lượng lao động).}
\item \cdkv{Cohesion}{"No longer... but..." (Cấu trúc nhấn mạnh), "Consequently" (Hệ quả).}
\item \cdkv{Comparison}{So sánh "Without this competence" với "being illiterate in the past" -> Làm nổi bật luận điểm "Fourth branch".}
\item \cdkv{Vocab}{Prerequisite (điều kiện tiên quyết), Competence (năng lực), Ubiquity (sự có mặt khắp nơi).}
\end{itemize}
\end{itemize}

\cdstep{3. Body Paragraph 2: Công cụ tiếp cận tri thức \textit{(Gateway to Knowledge)}}

"Furthermore, mastering computer skills is instrumental in acquiring knowledge in other disciplines, thereby reinforcing its role as a foundational branch of education. Just as reading is the key to unlocking literature, computer proficiency is the gateway to the vast resources available on the internet. Students equipped with strong IT skills can conduct research, access online libraries, and utilize educational software to enhance their understanding of math or science. Conversely, excluding this subject from the main curriculum would deprive students of a vital tool for self-education and lifelong learning, leaving them ill-equipped to navigate an increasingly digitalized society."
\begin{itemize}
\item \cdblue{\textbf{Phân tích:}}
\begin{itemize}
\item \cdkv{Argument}{Biến máy tính từ một "môn học" thành một "công cụ nền tảng" (foundational branch).}
\item \cdkv{Analogy (So sánh ẩn dụ)}{So sánh "Reading -> Literature" với "Computer -> Internet resources". Đây là cách lập luận rất thuyết phục.}
\item \cdkv{Vocab}{Instrumental in (đóng vai trò quan trọng trong), Deprive someone of something (tước đoạt của ai cái gì), Digitalized society (xã hội số hóa).}
\end{itemize}
\end{itemize}

\cdstep{4. Conclusion}

"In conclusion, I unequivocally support the view that computer skills should be standardized as the fourth essential branch of education alongside reading, writing, and mathematics. As technology continues to permeate every aspect of our lives, treating digital competency as a core subject is not merely a suggestion but an educational imperative."
\begin{itemize}
\item \cdblue{\textbf{Phân tích:}}
\begin{itemize}
\item \cdkv{Strong assertion}{"Unequivocally support" (Ủng hộ không chút nghi ngờ).}
\item \cdkv{Final thought}{"Not merely a suggestion but an educational imperative" (Không chỉ là gợi ý mà là mệnh lệnh giáo dục) -> Kết thúc mạnh mẽ.}
\item \cdkv{Vocab}{Permeate (thấm vào/lan tỏa vào), Imperative (điều cấp bách).}
\end{itemize}
\end{itemize}

\cdsection{PHẦN IV: TỪ VỰNG NÂNG CAO \textit{(ADVANCED VOCABULARY)}}

Học sinh cần tránh dùng lặp lại các từ đơn giản như "important", "good", "use computers". Hãy áp dụng các nhóm từ sau:

\cdkv{Nhóm 1}{Từ vựng về Tầm quan trọng \& Cốt lõi (Importance \& Core)}

Thay thế cho "Important", "Main subject"

Từ vựng / Cụm từ

Nghĩa tiếng Việt

Ví dụ (Context)

Indispensable / Imperative

Không thể thiếu / Cấp bách

Computer skills are indispensable in the modern era.

Cornerstone / Pillar

Nền tảng / Cột trụ

Maths is a cornerstone of science education.

Prerequisite

Điều kiện tiên quyết

IT skills are a prerequisite for most jobs.

Fundamental branch

Nhánh cơ bản/cốt lõi

It should be treated as a fundamental branch of learning.

Core curriculum

Chương trình học cốt lõi

Coding should be added to the core curriculum.

\cdkv{Nhóm 2}{Từ vựng về Công nghệ \& Kỹ năng (Tech \& Skills)}

Thay thế cho "Computer skills", "Using computers"

Từ vựng / Cụm từ

Nghĩa tiếng Việt

Ví dụ (Context)

Digital literacy

Sự hiểu biết về kỹ thuật số

Digital literacy is as important as reading.

Computer proficiency / Competency

Sự thành thạo máy tính

Students need to demonstrate computer proficiency.

Technological savvy

Sự am hiểu công nghệ

The younger generation is generally technologically savvy.

Digitalized society

Xã hội số hóa

We are living in a fully digitalized society.

Navigate the digital world

Xoay sở trong thế giới số

Schools must teach kids how to navigate the digital world.

\cdkv{Nhóm 3}{Cấu trúc câu ghi điểm (Grammar Structures)}

\cdstep{1. Cấu trúc so sánh nâng cao \textit{(Advanced Comparison)}:}
\begin{itemize}
\item Just as [A is important], so too is [B]. (Giống như A quan trọng, B cũng vậy).
\item \cdkv{Example}{Just as reading is crucial for literature, so too is coding for modern problem-solving.}
\end{itemize}

\cdstep{2. Mệnh đề quan hệ rút gọn \textit{(Reduced Relative Clause)}:}
\begin{itemize}
\item \cdkv{Example}{Students equipped with IT skills can learn faster. (Thay vì: Students who are equipped with...)}
\end{itemize}

\cdstep{3. Câu điều kiện loại ẩn \textit{(Implied Conditional - Without)}:}
\begin{itemize}
\item \cdkv{Example}{Without this competence, job seekers would face significant disadvantages. (Thay vì: If they do not have this competence...)}
\end{itemize}

\end{document}
